\documentclass[11pt]{ctexart}

\usepackage[margin=1in]{geometry}
\usepackage{amsmath,amssymb,amsthm,mathtools}
\usepackage{array}
\usepackage{booktabs}
\usepackage{enumitem}
\usepackage{hyperref}
\usepackage[nameinlink,noabbrev]{cleveref}
\usepackage{xcolor}
\usepackage{framed}

\hypersetup{
  colorlinks=true,
  linkcolor=blue!50!black,
  citecolor=blue!50!black,
  urlcolor=blue!50!black
}

\definecolor{codebg}{rgb}{0.95,0.95,0.95}

\title{TOOLS/SimplexNet 项目规划\\\large 面向强化学习的单纯形记忆网络工具包}
\author{Shuheng Zhang}
\date{\today}

\begin{document}

\maketitle

\begin{abstract}
本文稿描述 \texttt{TOOLS/SimplexNet} 项目的设计规划。该项目是一个面向强化学习的 Simplex Memory Network (SMN) 工具包，基于 \texttt{exp0414\_simplexNet} 的代码实现和 \texttt{writing/w0001-simplex-theory} 的理论基础。文稿分为三部分：其一，说明项目定位与目标用户；其二，描述两层类设计与两个子模块的组织架构；其三，给出实现优先级与 Phase 划分。本文稿旨在在代码实现前完成设计对齐，避免后续返工。
\end{abstract}

\tableofcontents

\section{项目定位}

\texttt{TOOLS/SimplexNet} 的核心目标是提供一个\textbf{开箱即用}的 SMN 强化学习工具包，使研究者和开发者能够快速将 SMN 集成到 RL 系统中。与 \texttt{exp0414\_simplexNet} 的"研究原型"定位不同，本项目的定位是\textbf{生产级工具包}。

目标用户包括：
\begin{itemize}[leftmargin=2em]
  \item 想要快速使用 SMN 进行 RL 实验的研究者；
  \item 需要将 SMN 集成到现有强化学习系统中的开发者；
  \item 想要评估 SMN vs 传统 MLP 的算法工程师。
\end{itemize}

核心价值体现在三个方面：
\begin{enumerate}[leftmargin=2em]
  \item \textbf{最小依赖}：仅需 PyTorch 和 Gymnasium，清晰接口；
  \item \textbf{可恢复训练}：完整的 checkpoint 生命周期管理；
  \item \textbf{向后兼容}：保持与 \texttt{exp0414} 的 API 互操作性。
\end{enumerate}

\section{核心类设计}

本项目对外暴露\textbf{两个核心类}，分别对应两个独立的 Python 文件。

\subsection{SMNmodule 类}

\texttt{SMNmodule} 继承自 \texttt{torch.nn.Module}，纯网络前向传播模块。

\begin{framed}
\begin{verbatim}
class SMNmodule(nn.Module):
    """Simplex Memory Network as a PyTorch module.

    Args:
        n: Simplex dimension (n=2 -> triangle, n=3 -> tetrahedron)
        m: Resolution (lattice points per edge)
        n_in: Input dimension
        n_out: Output dimension
        activation: Hidden node activation
        x_bounds: Per-channel input normalization bounds

    Example::
        # DQN Q-network (CartPole)
        q_net = SMNmodule(n=2, m=4, n_in=4, n_out=2)
    """

    def __init__(self, n=2, m=3, n_in=1, n_out=1,
                 activation='relu', x_bounds=None):
        super().__init__()
        # 初始化单纯形图、网络参数

    def forward(self, x: Tensor) -> Tensor:
        """Forward pass through the DAG."""
\end{verbatim}
\end{framed}

使用示例：
\begin{framed}
\begin{verbatim}
# DQN Q-network (CartPole)
q_net = SMNmodule(n=2, m=4, n_in=4, n_out=2)
output = q_net(state_tensor)
\end{verbatim}
\end{framed}

\subsection{SMN\_RL 类}

\texttt{SMN\_RL} 是 RL 功能封装类，组合了网络、算法、持久化和 GUI。

\begin{framed}
\begin{verbatim}
class SMN_RL:
    """Reinforcement Learning with SMN.

    Attributes:
        module (SMNmodule): The underlying SMN network
        env (gym.Env): Gymnasium environment
        algorithm (str): 'DQN', 'PPO', 'REINFORCE'
        gamma (float): Discount factor

    Methods:
        train(episodes, render): Train the agent
        test(episodes): Evaluate the agent
        plot(save_path): Visualize training curves
        save_checkpoint(path): Save training state
        load_checkpoint(path): Resume from checkpoint
        launch_gui(): Open PySide interactive window
    """
\end{verbatim}
\end{framed}

使用示例：
\begin{framed}
\begin{verbatim}
# 快速使用
env = gym.make('CartPole-v1')
agent = SMN_RL.create_default(env, algorithm='DQN')
agent.train(episodes=1000)
agent.launch_gui()  # 打开交互窗口
\end{verbatim}
\end{framed}

\subsection{类继承关系}

\begin{center}
\begin{tabular}{ll}
\toprule
类 & 继承/组合关系 \\
\midrule
\texttt{SMNmodule} & 继承 \texttt{nn.Module} \\
\texttt{SMN\_RL} & 组合 \texttt{SMNmodule} + \texttt{gym.Env} + 算法策略 \\
\bottomrule
\end{tabular}
\end{center}

\section{子模块设计}

项目包含两个子文件夹：\texttt{rl/} 和 \texttt{tools/}。

\subsection{rl/ --- 强化学习算法模块}

\begin{framed}
\begin{verbatim}
rl/
├── __init__.py
├── algorithms/
│   ├── dqn.py       # DQN 算法
│   ├── ppo.py       # PPO 算法 (Phase 2)
│   └── reinforce.py # REINFORCE 算法 (Phase 2)
├── mdp.py           # MDP 定义基类
└── env_wrapper.py   # 环境封装
\end{verbatim}
\end{framed}

DQN 算法接口示例：
\begin{framed}
\begin{verbatim}
from rl.algorithms import DQN

dqn = DQN(
    q_network=smn_module,
    obs_dim=4, act_dim=2,
    gamma=0.99, lr=1e-3
)

dqn.select_action(state, epsilon=0.1)
dqn.store_transition(state, action, reward, next_state, done)
loss = dqn.train_step()
\end{verbatim}
\end{framed}

\subsection{tools/ --- 工具模块}

\begin{framed}
\begin{verbatim}
tools/
├── __init__.py
├── checkpoint.py    # CheckpointManager
├── logger.py        # TrainingLogger
├── plot.py          # 可视化函数
└── gui.py           # PySide 交互窗口
\end{verbatim}
\end{framed}

\textbf{CheckpointManager} 负责数据持久化：
\begin{framed}
\begin{verbatim}
from tools import CheckpointManager

ckpt_mgr = CheckpointManager('./checkpoints')
ckpt_mgr.save(module=smn, optimizer=opt, episode=100)
checkpoint = ckpt_mgr.load_latest()
\end{verbatim}
\end{framed}

\textbf{TrainingLogger} 负责追加式日志（JSON Lines 格式）：
\begin{framed}
\begin{verbatim}
from tools import TrainingLogger

logger = TrainingLogger('./logs')
logger.log('train_start', config={...})
logger.log('epoch', episode=1, reward=100, loss=0.5)
\end{verbatim}
\end{framed}

\textbf{GUI 窗口} 提供 PySide 交互界面：
\begin{itemize}[leftmargin=2em]
  \item 参数调节滑块（n, m, lr, gamma, epsilon）
  \item 训练控制按钮（Start, Pause, Resume, Stop）
  \item 实时曲线显示（reward, loss, epsilon）
  \item Terminal log 面板
  \item Checkpoint 管理
\end{itemize}

\section{目录结构}

\begin{framed}
\begin{verbatim}
TOOLS/SimplexNet/
├── SMNmodule.py             # 主脚本 1: SMNmodule 类
├── SMN_RL.py                # 主脚本 2: SMN_RL 类
├── rl/                      # 强化学习模块
│   ├── algorithms/
│   ├── mdp.py
│   └── env_wrapper.py
├── tools/                   # 工具模块
│   ├── checkpoint.py
│   ├── logger.py
│   ├── plot.py
│   └── gui.py
├── __init__.py              # 导出：SMNmodule, SMN_RL
├── README.md
└── examples/
    ├── train_rl.py
    └── params.yaml
\end{verbatim}
\end{framed}

\section{实现优先级}

\subsection{Phase 1 (高优先级) --- 最小可用版本}

\begin{itemize}[leftmargin=2em]
  \item \texttt{SMNmodule.py} --- 网络模型
  \item \texttt{SMN\_RL.py} --- RL 封装类
  \item \texttt{rl/algorithms/dqn.py} --- DQN 算法
  \item \texttt{tools/checkpoint.py} --- CheckpointManager
  \item \texttt{tools/logger.py} --- TrainingLogger
  \item \texttt{tools/plot.py} --- 训练曲线可视化
  \item \texttt{\_\_init\_\_.py} + \texttt{README.md}
  \item \texttt{examples/train\_rl.py} --- RL 训练示例
\end{itemize}

\subsection{Phase 2 (中优先级) --- 增强功能}

\begin{itemize}[leftmargin=2em]
  \item \texttt{rl/algorithms/ppo.py} --- PPO 算法
  \item \texttt{rl/algorithms/reinforce.py} --- REINFORCE 算法
  \item \texttt{rl/mdp.py} --- MDP 定义基类
  \item \texttt{tools/gui.py} --- PySide 交互窗口
\end{itemize}

\subsection{Phase 3 (低优先级) --- 完善生态}

\begin{itemize}[leftmargin=2em]
  \item \texttt{tests/} --- 单元测试
  \item \texttt{examples/} --- 更多示例
  \item 文档网站
\end{itemize}

\section{与 exp0414 的对比}

\begin{center}
\begin{tabular}{p{0.25\linewidth} p{0.3\linewidth} p{0.3\linewidth}}
\toprule
方面 & exp0414 (现状) & SimplexNet (目标) \\
\midrule
定位 & 研究原型 + 实验验证 & 生产级工具包 \\
代码组织 & 扁平结构 (6 文件) & 模块化分层 \\
Checkpoint & 仅内存 best\_state & 完整磁盘管理 \\
日志 & 无 & 追加式训练日志 \\
GUI & 无 & PySide 交互窗口 \\
\bottomrule
\end{tabular}
\end{center}

\section{下一步行动}

计划已就绪，等待用户指令开始 Phase 1 实现。实现顺序：
\begin{enumerate}[leftmargin=2em]
  \item \texttt{SMNmodule.py}
  \item \texttt{rl/} 子文件夹
  \item \texttt{rl/algorithms/dqn.py}
  \item \texttt{tools/} 子文件夹
  \item \texttt{tools/checkpoint.py}
  \item \texttt{tools/logger.py}
  \item \texttt{tools/plot.py}
  \item \texttt{SMN\_RL.py}
  \item \texttt{\_\_init\_\_.py} + \texttt{README.md}
\end{enumerate}

\section{参考文档}

\begin{itemize}[leftmargin=2em]
  \item exp0414 工程说明文稿：\texttt{exp0414\_simplexNet/report\_tex/report.tex}
  \item SMN 理论文档：\texttt{writing/w0001-simplex-theory/report.tex}
  \item exp0414 代码：\texttt{exp0414\_simplexNet/src/}
  \item exp0420 RL 代码：\texttt{exp0420\_RLforSMN/src/}
\end{itemize}

\end{document}
